% ECE 442 / CSC 402 Project Proposal
% Don Nguyen

\documentclass[11pt]{article}

\usepackage[margin=1in]{geometry}
\usepackage{amsmath,amsfonts,amssymb}
\usepackage{graphicx}
\usepackage{hyperref}
\usepackage{url}
\usepackage{enumitem}

\title{
Agentic Relational Reasoning with Knowledge Graph Memory \\
for Enterprise Databases
}

\author{
  Don Nguyen \\
  CSC 402 / ECE 442 -- Network Science Analytics \\
  University of Rochester \\
  \texttt{knguy42@u.rochester.edu}
}

\begin{document}

\maketitle

\begin{abstract}
This proposal outlines a research-style project in network science analytics that builds a knowledge graph over enterprise databases to act as a context and memory layer for data querying agents.
The graph is constructed as a heterogeneous network over schemas, documentation, and other metadata from the Spider~2.0 and BIRD benchmark suites.
Using this graph as a structured memory, we aim to improve multi-hop SQL reasoning and complex query answering compared to strong large language model (LLM) baselines that do not leverage graph-based context.
\end{abstract}

\section{Problem Description}

Large language models (LLMs) have demonstrated impressive performance on text-to-SQL and conversational querying benchmarks, yet they still struggle on complex, multi-hop reasoning over realistic enterprise databases.
Benchmarks such as Spider~2.0 and BIRD highlight two key difficulties:
(1) very large, relationally rich schemas that are hard to load directly into a context window, and
(2) multi-turn interactions where the agent must remember prior steps and intermediate results.

Most current systems treat schema and documentation as flat text, using vector-based retrieval or heuristic pruning.
This approach largely ignores the intrinsic network structure of the database and its surrounding knowledge: tables, columns, foreign keys, business entities, and documentation snippets form a natural graph.
Without explicitly modeling this graph, agents may retrieve irrelevant tables, miss key relational paths, or fail to reuse past reasoning.

\textbf{Goal.}
The goal of this project is to build and analyze a \emph{knowledge graph memory layer} for enterprise-style databases, constructed from the schemas and documentation of Spider~2.0 and BIRD.
This graph will serve as a heterogeneous information network that a text-to-SQL agent can query to:
(i) prune schemas to relevant neighborhoods,
(ii) navigate multi-hop relational paths,
and (iii) maintain structured memory across interactions.
We will evaluate whether this network-science-based memory layer improves performance on BIRD and Spider~2.0 metrics compared to strong LLM baselines that operate without such a graph.

\section{Planned Approach}

\subsection{Network Construction}

We will model each benchmark database as a heterogeneous graph:
nodes will include tables, columns, higher-level entities, and documentation snippets;
edges will encode foreign-key relationships, table--column membership, semantic similarity links, and possibly \emph{usage} edges derived from query logs (when available).

\begin{itemize}[leftmargin=*]
  \item \textbf{Nodes:}
  \begin{itemize}
    \item Table nodes (one per table).
    \item Column nodes (one per column).
    \item Documentation nodes for relevant text snippets (e.g., schema descriptions, dataset documentation, FAQs).
  \end{itemize}
  \item \textbf{Edges:}
  \begin{itemize}
    \item Foreign key edges between columns/tables.
    \item Table--column membership edges.
    \item Semantic similarity edges (e.g., between columns and docs) using embedding-based similarity.
  \end{itemize}
\end{itemize}

The resulting network is a context graph or heterogeneous information network over which we can perform neighborhood queries, path searches, and other network-science operations.
We expect to start with a relatively small subset of Spider~2.0 and BIRD databases and scale up as compute and data access allow.

\subsection{Agent Architecture}

We plan to design an agentic text-to-SQL system where the LLM interacts explicitly with the knowledge graph:

\begin{enumerate}[leftmargin=*]
  \item \textbf{Schema and context retrieval.}
  Given a natural-language question (or multi-turn dialog), the agent first formulates a graph query (implicitly via an LLM prompt) to identify a subgraph neighborhood that is likely relevant to the question.
  This neighborhood includes a small set of tables, columns, and associated documentation.
  Graph queries will be implemented via \texttt{Neo4j} queries and/or graph search operations.
  \item \textbf{Reasoning and SQL generation.}
  The retrieved neighborhood is serialized into a compact textual or structured representation and provided to the LLM (e.g., via Nemotron NIM APIs).
  The LLM then reasons about the question, the schema neighborhood, and prior conversation turns to propose candidate SQL queries.
  \item \textbf{Execution feedback and memory update.}
  Candidate queries are executed (where possible) against the benchmark database.
  Results and execution feedback are used by the agent to refine its reasoning and can be written back to the graph as additional nodes/edges (e.g., ``successful path'' edges or ``frequently joined tables'').
  \item \textbf{Orchestration.}
  We plan to use LangGraph as a high-level orchestration framework, defining nodes for graph retrieval, LLM reasoning, SQL execution, and self-correction loops.
\end{enumerate}

\subsection{Baselines and Comparisons}

We will compare against baseline systems that \emph{do not} leverage a graph-structured memory:

\begin{itemize}[leftmargin=*]
  \item \textbf{Plain LLM:} A strong LLM (e.g., a Nemotron model accessed via NIM) prompted with the full or heuristically-pruned schema in text form, without graph reasoning.
  \item \textbf{Text-only RAG:} A vector-based retrieval system over schema and documentation text, where retrieved text chunks are given to the LLM, but no explicit graph or relational structure is exposed.
  \item \textbf{Existing benchmark baselines:} We will also use the official baselines and leaderboards for BIRD and Spider~2.0 when possible, as external reference points.
\end{itemize}

The core scientific question is whether modeling the schema and documentation as an explicit network yields measurable gains in complex, multi-hop reasoning performance over strong LLMs without such a network.

\section{Initial References}

We will build on the following initial set of references~\cite{yu2024spider2,li2024bird,sun2013survey,ji2022survey,schlichtkrull2018rgcn}:

\begin{itemize}[leftmargin=*]
  \item \textbf{Benchmarks.}
  \begin{itemize}
    \item Spider~2.0 benchmark paper and documentation (complex enterprise-scale text-to-SQL).
    \item BIRD benchmark for multi-turn, conversational, and agentic SQL reasoning.
  \end{itemize}
  \item \textbf{Network science and graph analytics.}
  \begin{itemize}
    \item Standard references on heterogeneous information networks, graph representation learning, and network analytics (e.g., textbooks and survey papers from the course resources).
  \end{itemize}
  \item \textbf{Graph-based and agentic SQL systems.}
  \begin{itemize}
    \item Recent work on graph-based schema representations and relational foundation models (e.g., systems that treat databases as graphs for retrieval and reasoning).
    \item Papers on agentic workflows for text-to-SQL, tool-augmented LLMs, and memory-augmented agents.
  \end{itemize}
\end{itemize}

As part of the project, we will expand this list into a more thorough literature review, covering both classical network science methods and modern graph neural networks / relational models applied to database and knowledge graph reasoning.

\section{Software, Tools, and Infrastructure}

We plan to use the following software stack:

\begin{itemize}[leftmargin=*]
  \item \textbf{LLM APIs:} NVIDIA Nemotron NIM for primary model inference, with potential use of additional LLM APIs for comparison.
  \item \textbf{Graph layer:} Neo4j as the primary graph database to store and query the heterogeneous knowledge graph.
  \item \textbf{Orchestration:} LangGraph for defining the agent workflow (graph retrieval, LLM calls, SQL execution, self-correction).
  \item \textbf{General programming environment:} Python-based stack for data preprocessing, graph construction, and evaluation scripts.
\end{itemize}

Because full Spider~2.0 evaluation requires access to cloud data warehouses such as Snowflake or Databricks, we will start with subsets or simplified setups and escalate to more realistic deployments only if resources permit.

\section{Datasets and Networks}

The primary datasets and networks are:

\begin{itemize}[leftmargin=*]
  \item \textbf{BIRD.}
  We will use the BIRD benchmark for multi-turn, conversational SQL reasoning.
  From BIRD, we will extract database schemas, foreign key relationships, and associated documentation to construct heterogeneous networks.
  BIRD will serve as the \emph{primary} evaluation benchmark if Spider~2.0 access is limited.
  \item \textbf{Spider~2.0.}
  We will use Spider~2.0, focusing on the complex enterprise schemas that require cloud data warehouses.
  Here the emphasis is on modeling large, relationally rich networks and testing how well the graph memory scales.
  Full evaluation on Spider~2.0 is considered a stretch goal, dependent on access to the necessary cloud infrastructure.
\end{itemize}

In all cases, the networks we analyze are derived from the schemas and documentation: they are not arbitrary graphs, but structured representations of real-world relational databases and their knowledge context.

\section{Expected Outcomes and Success Criteria}

\subsection{Minimum Viable Success}

The minimum viable outcome is an agentic text-to-SQL system that:

\begin{itemize}[leftmargin=*]
  \item Implements a knowledge graph memory over at least a subset of BIRD and Spider~2.0 schemas and documentation.
  \item Demonstrates measurable improvements in benchmark scores on BIRD (e.g., execution accuracy or success rate) compared to:
  (i) a plain LLM baseline and
  (ii) a text-only RAG baseline.
\end{itemize}

\subsection{Quantitative Criteria}

We will primarily use the official metrics defined by BIRD and Spider~2.0, such as:

\begin{itemize}[leftmargin=*]
  \item \textbf{Execution accuracy:} whether the generated SQL returns the correct result.
  \item \textbf{Exact-match or partial-match metrics:} string-level or structure-level comparison to reference SQL when available.
  \item \textbf{Multi-turn success rate:} for conversational tasks, the fraction of dialogues in which the agent successfully answers all required questions.
\end{itemize}

We will also perform ablation studies comparing:

\begin{enumerate}[leftmargin=*]
  \item No graph (plain LLM).
  \item Text-only retrieval over schema/docs.
  \item Full knowledge graph memory (proposed method).
\end{enumerate}

\section{Additional Notes, Risks, and Stretch Goals}

\subsection{Risks and Constraints}

A key practical risk is the infrastructure required to fully evaluate on Spider~2.0, which relies on cloud data warehouses such as Snowflake or Databricks.
If access to these resources is limited, we may only be able to test smaller subsets of Spider~2.0 or use local emulations.

Another constraint is time and compute budget for running large numbers of LLM calls and SQL executions, especially for multi-turn conversational evaluation.
We will mitigate this by starting with small-scale experiments and carefully budgeting benchmark runs.

\subsection{Stretch Goals}

If resources and time permit, we will pursue the following stretch goals:

\begin{itemize}[leftmargin=*]
  \item Full-scale evaluation on Spider~2.0 with realistic cloud warehouse setups.
  \item More advanced graph-learning methods over the knowledge graph (e.g., graph neural networks for better node/edge embeddings).
  \item Incorporating dynamic memory updates into the graph based on agent interactions (e.g., recording successful query patterns or typical join paths).
\end{itemize}

Even if these stretch goals are not fully realized, the project will deliver a concrete network-based agent architecture, an implemented knowledge graph memory layer, and empirical evidence about the value of graph-structured context for enterprise text-to-SQL reasoning.

\begin{thebibliography}{9}

\bibitem{yu2024spider2}
Tao Yu, et al.
\newblock Spider 2.0: A multi-database benchmark for realistic, large-scale text-to-SQL evaluation.
\newblock \emph{ArXiv preprint arXiv:2407.XXXX}, 2024.

\bibitem{li2024bird}
Haoyang Li, et al.
\newblock BIRD: A benchmark for intelligent reasoning over databases.
\newblock \emph{ArXiv preprint arXiv:2401.XXXX}, 2024.

\bibitem{sun2013survey}
Yizhou Sun and Jiawei Han.
\newblock Mining heterogeneous information networks: a structural analysis approach.
\newblock \emph{ACM SIGKDD Explorations Newsletter}, 14(2):20--28, 2013.

\bibitem{ji2022survey}
Shuai Ji, et al.
\newblock A survey on knowledge graphs: representation, acquisition, and applications.
\newblock \emph{IEEE Transactions on Neural Networks and Learning Systems}, 33(2):494--514, 2022.

\bibitem{schlichtkrull2018rgcn}
Michael Schlichtkrull, et al.
\newblock Modeling relational data with graph convolutional networks.
\newblock In \emph{Proceedings of the European Semantic Web Conference (ESWC)}, 2018.

\end{thebibliography}

\end{document}

